\section{Proofs for Types}

\begin{example}
  We generalize the example from the introduction to a
  parametric stack protocol.
\begin{lstlisting}
protocol Stack a = Pop -a | Push a (Stack a) (Stack a)
\end{lstlisting}
  We set  $ \Delta = \types{Neg} : \kindp \to \kindp, \types{a} :
  \kindp$ and obtain
  the straightforward derivation
  \begin{mathpar}
    \inferrule*{
      \inferrule*
      {\inferrule*
      {\typeSynth{ \types{a} }{\kindp}}
      {\typeAgainst{ \types{a} }{\kindp}}}
      {\typeAgainst{ \types{-a} }{\kindp}} \\
      \inferrule*
      {\typeSynth{ \types{a} }{\kindp}}
      {\typeAgainst{ \types{a} }{\kindp}} \\
      \inferrule*
      {\inferrule*
        {\types{Stack} : {\kindp} \to \kindp \in \Delta \\
          \inferrule*
          {\typeSynth{\types{a}}{\kindp}}
          {\typeAgainst{\types{a}}{\kindp}}}
        {\typeSynth{ \types{Stack~a}}{\kindp}}}
      {\typeAgainst{ \types{Stack~a}}{\kindp}}
    }
    {
      \typeSynth[]{ \types{Stack} }{\kindp \to \kindp}
    }
  \end{mathpar}
  The three subderivations correspond to the arguments of the protocol
  constructors. We conclude that the definition of the protocol 
  \lstinline+Stack a+ is well-formed.
\end{example}

\begin{lemma}[Properties of the materialization and directional operators]\
  \label{lem:ops-dual}
  \begin{enumerate}
    \item\label{it:dual-norm} $\isConv{\TDual{(\tosession TS)}}{\tosession{\polarizedparam{-}{T}}{\TDual{S}}}\kinds$
    \item\label{it:plus-minus} $\polarizedparam{-}{\polarizedparam{+}{T}} = \polarizedparam{-}{T}$
    \item\label{it:minus-minus} $\polarizedparam{-}{\polarizedparam{-}{T}} = \polarizedparam{+}{T}$
  \end{enumerate}
\end{lemma}
\begin{proof}\
  \begin{enumerate}
  \item By case analysis on the definition of the
    $\tosessionop$ operator (\cref{fig:type-polarized-operator}).

    Case $\typec T = \pneg U $.
    %$\isConv{\TDual{(\tosession {-T}S)}}{\tosession{\polarizedparam{-}{-T}}{\TDual{S}}}\kinds$.
    Since $\tosession {\pneg U}S = \TIn{U}{S}$, rule \rulenametconvDualInM gives
    $\isConv{\TDual{(\tosession {\pneg U}S)}}{\TOut{U}{\TDual S}}\kinds$. Observing that
    $\polarizedparam{-}{-U}=\polarizedparam{+}{U}=\typec U$ we conclude that
    $\tosession{\polarizedparam{-}{\pneg U}}{\TDual{S}}=\TOut{U}{\TDual S}$.

    Case $\typec T \neq \pneg U $
    %$\isConv{\TDual{(\tosession {T}S)}}{\tosession{\polarizedparam{-}{T}}{\TDual{S}}}\kinds$.
    Since $\tosession {T}S = \TOut{T}{S}$, rule \rulenametconvDualOutM gives
    $\isConv{\TDual{(\tosession {T}S)}}{\TIn{T}{\TDual S}}\kinds$. Conclude observing
    that $\polarizedparam{-}{T}=\typec{-T}$.
    \item The proof follows by case analysis on the definition of the directional operators.

    Case $\typec T = \pneg U $. We have:
    $\polarizedparam{-}{\polarizedparam{+}{\pneg U}} = \polarizedparam{-}{\polarizedparam{-}{U}}
    = \polarizedparam{-}{\pneg{U}}$.

    Case $\typec T \neq \pneg U $. We have:
    $\polarizedparam{-}{\polarizedparam{+}{T}} = \polarizedparam{-}{T}$.

    \item The proof follows by case analysis on the definition of the directional operators.

    Case $\typec T = \pneg U $. We have:
    $\polarizedparam{-}{\polarizedparam{-}{\pneg U}} = \polarizedparam{-}{\polarizedparam{+}{U}}
    = \polarizedparam{-}{U} = \pneg U = \polarizedparam{+}{\pneg U}$.

    Case $\typec T \neq \pneg U $. We have:
    $\polarizedparam{-}{\polarizedparam{-}{T}} = \polarizedparam{-}{\pneg T}= \polarizedparam{+}{T}$.
  \end{enumerate}
\end{proof}

\begin{lemma}[Agreement for type conversion]
  \label{it:agreement-type-conv}
  If $\isConv T U \kind$, then $\typeAgainst T \kind$ and
  $\typeAgainst U \kind$.
\end{lemma}
\begin{proof}
  By rule induction on the hypothesis.
\end{proof}


We start by specifying the syntax of types in normal form. Normal forms for well-formed types
are $\typec{Q}$ as defined by the following grammar:
%
\begin{align*}
  \typec{Q} &\grmeq \typec{R} \grmor \pneg{R}
  \\
  \typec{R} &\grmeq
  % functional types
  \TUnit \grmor \TFun[]{R}{R} \grmor \TPair{R}{R} \grmor
  \TForall \alpha \kind {R} \grmor \typec\alpha
  % session types
  \grmor \TIn{R}{R} \grmor \TOut{R}{R}
  \grmor \TEndW \grmor \TEndT  \grmor
  \TDual{\alpha}
  % protocol and data types
  \grmor \TPApp \proto {\overline{Q}}
\end{align*}


\begin{proposition}[Kind preservation]
  ~\label{lemma:kind-preservation}
  \begin{enumerate}
  \item For all $\typeSynth T\kind$,  $\isConv{\Normal[+] T}T\kind$.
  \item For all $\typeSynth T\kinds$,  $\isConv{\Normal[-] T}{\TDual T}\kinds$.
  \end{enumerate}
\end{proposition}
\begin{proof}
  See Agda development \texttt{nf-sound+} and \texttt{nf-sound-}.
\end{proof}

\begin{proof}[Proof of \cref{thm:soundness-type-equiv} (Soundness)]
  \Cref{it:soundness-type-equiv-t}. Suppose that $\Normal[+] T =_\alpha \Normal[+] U$. From
  \cref{lemma:kind-preservation}, we have  $\isConv{\Normal[+]
    T}T\kind$ and  $\isConv{\Normal[+] U}U\kind$. Conclude by symmetry
  and transitivity of conversion as type equality is in $\equiv$.
  % 
  The proof for \cref{it:soundness-type-equiv-s} is analogous.
\end{proof}

\begin{proof}[Proof of \cref{thm:completeness-type-equiv} (Completeness)]
  See Agda development \texttt{nf-complete} and \texttt{nf-complete-}.
\end{proof}

To obtain a complexity bound for the computation of a normal form,
we exploit its syntactic form.
\begin{lemma}
  \label{lemma:syntax-of-type-normal-forms}
  If $\Normal[\pm]T = \typec{T'}$, then $\typec{T'}$ is described by $\typec{Q}$.
\end{lemma}
\begin{proof}
  See Agda development \texttt{nf-normal}.
\end{proof}
\begin{proposition}
  For each $\typeSynth T\kind$,
  the worst case running time of $\Normal[\pm]T$ is in $O(|\typec{T}|)$.
\end{proposition}
\begin{proof}
  A run of $\Normal T$ visits each of the $|\typec{T}|$ nodes in the input type
  once. At each node, it either reconstructs the type from the normal forms of
  the sub-terms in $O(1)$ time, or (in the transmission rules) it applies
  polarity correction to normal forms. By
  Lemma~\ref{lemma:syntax-of-type-normal-forms}, a normal form starts with at
  most one negative polarity, so correction takes at most two steps, at cost
  $O(1)$.
\end{proof}

\begin{proof}[Proof of \cref{thm:complexity} (Time complexity of type equivalence)]
  A corollary of \cref{lemma:syntax-of-type-normal-forms}.
\end{proof}

%%% Local Variables:
%%% mode: latex
%%% TeX-master: "main"
%%% End:
